\input{../tex-inputs/latex-leseansicht-vorspann}

\section[Arthur Schnitzler, Olga und Elisabeth Gussmann an Gustav Schwarzkopf, 12. 7. 1901]{L04427 Arthur Schnitzler, Olga und Elisabeth Gussmann an Gustav Schwarzkopf, 12. 7. 1901}
\nopagebreak\mylabel{L04427v}
\rehead{ }\normalsize\beginnumbering\briefempfaengerindex{Schwarzkopf, Gustav@\textsc{Schwarzkopf, Gustav}!zzzSteinrück, Elisabeth@\emph{von Elisabeth Steinrück}!1901-07-122@{12. 7. 1901}|(be}\briefempfaengerindex{Schwarzkopf, Gustav@\textsc{Schwarzkopf, Gustav}!zzzSchnitzler, Olga@\emph{von Olga Schnitzler}!1901-07-122@{12. 7. 1901}|(be}\briefempfaengerindex{Schwarzkopf, Gustav@\textsc{Schwarzkopf, Gustav}!zzzSchnitzler, Arthur@\emph{von Arthur Schnitzler}!1901-07-122@{12. 7. 1901}|(be}
\toendnotes[C]{\smallbreak\pagebreak[2]}
\correspDesc{Versand  durch Arthur Schnitzler, Olga Schnitzler, Elisabeth Steinrück am 12. 7. 1901 in Innsbruck
\newline{}Übermittlung  am 13. 7. 1901 in Innsbruck
\newline{}Zustellung  am 14. 7. 1901 in Wien
\newline{}Erhalt  durch Gustav Schwarzkopf am 14. 7. 1901 in Wien}\toendnotes[C]{\smallbreak}
\Standort{DLA, A:Schnitzler, HS.1985.1.1897.}
\physDesc{Bildpostkarte, 184 Zeichen
\newline{}Handschrift Arthur Schnitzler: Bleistift, deutsche Kurrent
\newline{}Handschrift Elisabeth Steinrück: Bleistift, lateinische Kurrent
\newline{}Handschrift Olga Schnitzler: Bleistift, deutsche Kurrent
\newline{}Versand: 1) Stempel: »\nobreak{}\oindex{Innsbruck@\textbf{Innsbruck}|pwk}Inn\textcolor{gray}{sbruck}, 13 7 0{[}1{]}\nobreak{}«.   2) Stempel: »\nobreak{}\oindex{I., Innere Stadt@\textbf{I., Innere Stadt}|pwk}Wien 1/1 1, 14 7 \textcolor{gray}{01}, \textcolor{gray}{9–10}½V., Bestellt\nobreak{}«. }\toendnotes[C]{\smallbreak}\pstart{}{\pb}Herrn Guſtav Schwarzkopf\pend{}\pstart{}\textsc{Wien}\oindex{Wien@\textbf{Wien}|pw}\pend{}\pstart{}\textsc{I. Tiefer Graben 23}\oindex{Wien@\textbf{Wien}!I., Innere Stadt@\textbf{I., Innere Stadt}!Tiefer Graben 23@\textbf{Tiefer Graben 23}|pw}.\pend{}{\bigskip}
\pstart
           \centering{}{\pb}\textcolor{gray}{\textbf{INNSBRUCK\oindex{Innsbruck@\textbf{Innsbruck}|pw} VON NORD.}}\pend
           \vspace{1em}
\pstart
           {\pb}12/7. 901.\pend
           \vspace{0.5em}
\pstart
           Dank für den lieben \label{K_L04427-1v}\edtext{Brief}{\lemma{\textnormal{\emph{Brief}}}\Cendnote{\textnormal{XXXX Auszeichnungsfehler: Dokument L04308 nicht gefunden.}}}\label{K_L04427-1}.\pend
           
\pstart
           Reiſen eben nach \textsc{Vahrn\oindex{Vahrn@\textbf{Vahrn}|pw}},\pend
           \pstart herzlich{[}e{]} Grüße. \spacefill\mbox{A}\pend{}\selectlanguage{ngerman}\vspace{1em}
\pstart
           \noindent{}{[}hs. Steinrück:{]} Servus,\pend
           \pstart \spacefill\mbox{Lisl}\pend{}\selectlanguage{ngerman}\vspace{1em}
\pstart
           \noindent{}{[}hs. O. Schnitzler:{]} Beſte Grüße!\pend
           \pstart \spacefill\mbox{Olga}\pend{}
\pstart
           \noindent{}\label{T_L04427-1v}\edtext{Wenn auch verdreht u. in den Wolken.}{\lemma{\textnormal{\emph{Wenn … Wolken.}}}\Cendnote{\textnormal{Ihr Text ist oberhalb und verkehrt zum Bildmotiv geschrieben, so dass er über die abgebildeten Wolken läuft.}}}\label{T_L04427-1}\pend
           \selectlanguage{ngerman}\endnumbering\briefempfaengerindex{Schwarzkopf, Gustav@\textsc{Schwarzkopf, Gustav}!zzzSteinrück, Elisabeth@\emph{von Elisabeth Steinrück}!1901-07-122@{12. 7. 1901}|)be}\briefempfaengerindex{Schwarzkopf, Gustav@\textsc{Schwarzkopf, Gustav}!zzzSchnitzler, Olga@\emph{von Olga Schnitzler}!1901-07-122@{12. 7. 1901}|)be}\briefempfaengerindex{Schwarzkopf, Gustav@\textsc{Schwarzkopf, Gustav}!zzzSchnitzler, Arthur@\emph{von Arthur Schnitzler}!1901-07-122@{12. 7. 1901}|)be}\mylabel{L04427h}
\begin{anhang}
\end{anhang}\newcommand{\dateiname}{L04427}\newcommand{\titel}{Arthur Schnitzler, Olga und Elisabeth Gussmann an Gustav Schwarzkopf, 12. 7. 1901}\newcommand{\editorInnen}{Herausgegeben von Jahnke, SelmaMüller, Martin Anton}\input{../tex-inputs/latex-leseansicht-abspann}
